% Converted from README.md
\section{Simboličko izvršavanje (alat KLEE)}\label{simboliux10dko-izvrux161avanje-alat-klee}

\href{http://klee.github.io/}{KLEE} je javno dostupan alat koji se distribuira pod licencom NCSA otvorenog koda Univerziteta Ilinois. Služi za simboličko izvršavanje programa i za automatsko generisanje test primera. KLEE vrši analizu nad LLVM međureprezentacijom i koristi SMT rešavač \href{https://stp.github.io/}{STP} za proveravanje uslova ispravnosti koje generiše. KLEE koristi nekoliko optimizacija i heuristika za poboljšavanje pokrivenosti koda prilikom simboličkog izvršavanja. KLEE se koristi i kao sastavni deo raznih platformi za razvijanje novih alata za analizu programa. KLEE ima dva cilja:

\begin{enumerate}
\def\labelenumi{\arabic{enumi}.}
\tightlist
\item
  da pokrije svaku liniju izvornog koda u programu
\item
  da detektuje svaku opasnu operaciju ako postoji ijedna ulazna vrednost koja može da prouzrokuje grešku
\end{enumerate}

Ukoliko prilikom generisanja test primera za neku od putanja, KLEE pronađe test primer koji izaziva grešku, onda će pored tog test primera biti generisane još dve datoteke:

\begin{itemize}
\tightlist
\item
  \texttt{test\textless{}N\textgreater{}.pc} - uslovi ograničenja koji važe za tu putanju i koji su prosleđeni solveru u \texttt{KQuery} sintaksi.
\item
  \texttt{test\textless{}N\textgreater{}.\textless{}TYPE\textgreater{}.err} - dodatne informacije o grešci
\end{itemize}

\texttt{\textless{}TYPE\textgreater{}} označava tip detektovane greške. Neke od grešaka koje KLEE može da detektuje su:

\begin{itemize}
\tightlist
\item
  \textbf{ptr}: Čitanje ili pisanje u nevalidnu memorijsku lokaciju;
\item
  \textbf{free}: Dupli ili nekorektan poziv funkcije \texttt{free()};
\item
  \textbf{abort}: Program je pozvao funkciju \texttt{abort()};
\item
  \textbf{assert}: Uslov proveravan funkcijom \texttt{assert} nije tačan;
\item
  \textbf{div}: Deljenje nulom ili računanje po modulu 0;
\item
  \textbf{user}: Problem je u vezi sa ulazom (neispravan poziv \texttt{klee\ intrinsic}) ili pogrešan način upotrebe KLEE-a.
\item
  \textbf{exec}: Dogodio se problem koji je KLEE sprečio u izvršavanju programa, npr. nepoznata instrukcija, poziv preko neispravnog pokazivača na funkciju ili inline funkcije.
\item
  \textbf{model}: KLEE nije mogao da ispuni kompletnu preciznost i ispituje samo deo svih stanja programa. Na primer, simboličke veličine prosleđene \texttt{malloc} funkciji nisu trenutno podržane i u takvim situacijama KLEE konkretizuje taj argument.
\end{itemize}

\subimport{01_intro/}{index.tex}

\subimport{02_err_semantic/}{index.tex}

\subimport{03_err_pointer/}{index.tex}

\subimport{04_libc/}{index.tex}

\subimport{05_heuristics/}{index.tex}
